
\clearpage
\chapter{מסגרות תיאום מודרניות}

\section{\textenglish{AutoGen}: שיחות רב-סוכן}

\textenglish{AutoGen} \cite{wu2023autogen} הנה מסגרת שיחות בין סוכנים: \textenglish{AssistantAgent} ו-\textenglish{UserProxyAgent} מנהלים שיח מובנה שבו הסוכן מציע קוד, ה-\textenglish{Proxy} מריץ אותו ומחזיר תוצאות. המסגרת תומכת בדפוסי שיחה מורכבים: \textenglish{group chat}, \textenglish{nested chat}, ו-\textenglish{sequential chat}.

\section{\textenglish{MetaGPT}: תפקידים מקצועיים}

\textenglish{MetaGPT} \cite{hong2023metagpt} מבוסס על מטאפורת "חברת תוכנה": לכל סוכן תפקיד מוגדר (מנהל מוצר, אדריכל, מתכנת, בודק). הניתוב מסתמך על \textenglish{Standard Operating Procedures (SOPs)} קשיחים, מה שמפחית רעש בבחירת כלים.

\section{\textenglish{CrewAI}: תפקידים גמישים}

\textenglish{CrewAI} מציע פשרה בין קשיחות \textenglish{MetaGPT} לגמישות \textenglish{AutoGen}: תצורת \textenglish{YAML} מגדירה כישורים, סוכנים ומשימות, ומנוע \textenglish{Process.HIERARCHICAL} מנתב עבודה דרך מנהל. מחקר זה \cite{nakash2024skillsieve} מוסיף שכבת אבטחה (\textenglish{SkillSieve}) מעל מנגנון הניתוב.

\begin{figure}[H]
\centering
\includegraphics[width=0.85\textwidth]{assets/architecture.png}
\caption{\textenglish{Timeline: Evolution of Multi-Tool Orchestration Frameworks (Python-generated)}}
\label{fig:arch}
\end{figure}
